\documentclass[12pt,a4paper]{article}
\usepackage{amsmath,amssymb,amsthm}
\usepackage[margin=1in]{geometry}
\usepackage{hyperref}
\usepackage{booktabs}
\hypersetup{colorlinks=false}

\newtheorem{theorem}{Theorem}
\newtheorem{proposition}[theorem]{Proposition}
\newtheorem{corollary}[theorem]{Corollary}

\title{\textbf{$K$-deviation as a vacuum diagnostic\\
in accreting Kerr black holes:\\
a self-consistency prediction for XRISM}}

\author{Y.\,Y.\,N.\ Li\\
\small Independent Researcher}
\date{}

\begin{document}
\maketitle

\begin{abstract}
The $K=1$ self-consistency framework~\cite{K1note} establishes
that the functional $K:=\Sigma\,\Box_{\rm null}\ln\sqrt\Sigma$
equals unity identically for the Kerr vacuum, and that
$K=1\Leftrightarrow R_{\mu\nu}\ell^\mu\ell^\nu=0$ for the
principal null congruence.
For a general perfect fluid with energy density $\rho$ and
pressure $p$, the Einstein equations give
\begin{equation}\label{eq:deltaK}
  \delta K := K - 1 = -4\pi\Sigma(\rho+p)(u_\mu\ell^\mu)^2,
\end{equation}
where $\Sigma=r^2+a^2\cos^2\!\theta$ and $\ell^\mu$ is the
Kinnersley outgoing null vector.
This is exact.
For accreting black holes, $\delta K<0$: matter reduces $K$
below unity by an amount proportional to $\Sigma(\rho+p)$.

We derive three observable signatures of $\delta K$ in the
Fe~K$\alpha$ reflection spectrum:
(i)~a systematic redshift $\delta E\approx|\delta K|\times 4.5$\,keV
of the iron line centroid;
(ii)~a linear correlation $\delta E(t)\propto L_X(t)$ between
the line shift and X-ray luminosity, which standard Kerr
reflection does not predict;
(iii)~an inclination dependence $\delta K\propto\Sigma(\theta)$
with amplitude $(r^2+a^2)/r^2-1$ at the pole versus equator.

For Cyg~X-1 ($a\approx0.98M$, near-ISCO accretion density
$\rho_0\sim10^{-3}/M^2$), we predict
$\delta K\approx-46\%$ and a line centroid shift
$\delta E\approx 2$\,keV---approximately $280\sigma$ above
the XRISM/Resolve energy resolution of 7\,eV.
For M87* ($\rho_0\sim10^{-4}/M^2$),
$\delta K\approx-4.6\%$ and $\delta E\approx 0.2$\,keV.
We propose that the luminosity correlation (signature~ii)
provides the cleanest separation of $\delta K$ from standard
spin-measurement systematics, and invite comparison with
existing XRISM data.
\end{abstract}

\medskip
\noindent\textbf{Keywords:} Kerr black holes $\cdot$
self-consistency $\cdot$ iron line $\cdot$ XRISM $\cdot$
accretion $\cdot$ vacuum diagnostic

%===================================================================
\section{Introduction}\label{sec:intro}

The Fe~K$\alpha$ reflection line at 6.4\,keV is a primary
diagnostic of black hole spin in X-ray binaries and active
galactic nuclei~\cite{Fabian2000,Reynolds2014}.
Standard analyses fit the observed line profile with
relativistic reflection models ({\tt RELXILL}~\cite{Garcia2014})
that assume the spacetime is \emph{exactly} the Kerr vacuum.
This assumption ignores the accreting matter itself, which
modifies the spacetime geometry at the level
$\delta g_{\mu\nu}\sim 8\pi G\rho r^2/c^4$.

The $K=1$ self-consistency framework~\cite{K1note} provides
a precise language for quantifying the departure from Kerr vacuum.
The functional
\begin{equation}
  K := \Sigma\,\Box_{\rm null}\ln\sqrt{\Sigma},\quad
  \Box_{\rm null} := \partial_r^2 + \frac{2r}{\Sigma}\partial_r,
\end{equation}
satisfies $K=1$ identically for the Kerr metric~\cite{K1note},
and---via the Raychaudhuri equation---$K=1$ is equivalent to
the vacuum condition $R_{\mu\nu}\ell^\mu\ell^\nu=0$ along the
principal outgoing null congruence.

When matter is present, $K$ departs from unity by
equation~\eqref{eq:deltaK}.
This departure has three observable consequences in the
Fe~K$\alpha$ line that are distinct from, and independent of,
standard spin-measurement effects.

%===================================================================
\section{The $\delta K$ formula}\label{sec:deltaK}

\begin{proposition}[Matter correction to $K$]\label{prop:deltaK}
For the Kerr metric with a perfect fluid stress-energy tensor
$T_{\mu\nu}=(\rho+p)u_\mu u_\nu + p\,g_{\mu\nu}$,
the $K$-functional satisfies
\begin{equation}\label{eq:deltaK_full}
  K = 1 - 4\pi\Sigma(\rho+p)(u_\mu\ell^\mu)^2.
\end{equation}
\end{proposition}

\begin{proof}
From the Proposition of~\cite{K1note},
$K = 1 - \tfrac{1}{2}\Sigma\,R_{\mu\nu}\ell^\mu\ell^\nu$.
The Einstein equations give
$R_{\mu\nu}\ell^\mu\ell^\nu = 8\pi T_{\mu\nu}\ell^\mu\ell^\nu$.
For a perfect fluid,
$T_{\mu\nu}\ell^\mu\ell^\nu = (\rho+p)(u_\mu\ell^\mu)^2 + p\,g_{\mu\nu}\ell^\mu\ell^\nu
= (\rho+p)(u_\mu\ell^\mu)^2$,
since $\ell^\mu$ is null ($g_{\mu\nu}\ell^\mu\ell^\nu=0$).
Substituting gives~\eqref{eq:deltaK_full}.\qed
\end{proof}

\noindent
For a pressureless dust ($p=0$) comoving with the fluid
($u_\mu\ell^\mu=1$):
\begin{equation}\label{eq:dust}
  \delta K = -4\pi\Sigma\rho_0 = -4\pi(r^2+a^2\cos^2\!\theta)\,\rho_0.
\end{equation}
For radiation-dominated accretion ($p=\rho/3$):
$\delta K = -\tfrac{16\pi}{3}\Sigma\rho$.

Three properties of~\eqref{eq:dust} are immediate:
\begin{enumerate}
\item \emph{Sign}: $\delta K<0$ for any matter ($\rho>0$).
  Matter always pushes $K$ below unity.
\item \emph{Scaling}: $|\delta K|\propto\Sigma\rho_0\propto\rho_0 r^2$
  at large $r$; near the ISCO, $\Sigma$ is $O(M^2)$.
\item \emph{Angular dependence}: $|\delta K|$ is larger at the
  poles ($\theta=0$, $\Sigma=r^2+a^2$) than at the equator
  ($\theta=\pi/2$, $\Sigma=r^2$) by a factor $(r^2+a^2)/r^2$.
\end{enumerate}

%===================================================================
\section{Observable signatures}\label{sec:obs}

\subsection{Signature~I: Iron line centroid shift}

The Fe~K$\alpha$ line centroid energy observed at infinity is
\begin{equation}
  E_{\rm obs} = E_0\,g(r,a,i),
\end{equation}
where $E_0=6.4$\,keV and $g$ is the redshift factor.
Near the ISCO, $g\approx0.7$ for typical inclinations.

The $\delta K$ correction modifies the effective gravitational
potential, shifting the line centroid by
\begin{equation}\label{eq:deltaE}
  \delta E \approx E_0\,g\,|\delta K|.
\end{equation}
This is a systematic redshift: accreting matter shifts the
line to lower energies relative to the pure Kerr prediction.

\subsection{Signature~II: Luminosity correlation}

Since $\rho_0\propto\dot{M}\propto L_X$, equation~\eqref{eq:dust}
gives
\begin{equation}\label{eq:correlation}
  \delta E(t) \propto L_X(t).
\end{equation}
The line centroid shift tracks the X-ray luminosity.
Standard Kerr reflection effects (gravitational redshift, Doppler
broadening) do not vary with luminosity for a fixed geometry.
This correlation is therefore a \emph{clean discriminant} between
$\delta K$ and spin-measurement systematics.

Table~\ref{tab:lum} gives predicted $\delta K$ and $\delta E$
for Cyg~X-1 ($a=0.98M$, $r=6M$, equatorial) as a function of
luminosity.

\begin{table}[h]
\centering
\renewcommand{\arraystretch}{1.3}
\begin{tabular}{@{}cccc@{}}
\toprule
$L/L_{\rm Edd}$ & $\rho_0\,[M^{-2}]$ & $\delta K$ & $\delta E$\,[eV] \\
\midrule
0.01 & $10^{-5}$ & $-0.005$ & 20 \\
0.05 & $5\times10^{-5}$ & $-0.023$ & 101 \\
0.10 & $10^{-4}$ & $-0.045$ & 203 \\
0.30 & $3\times10^{-4}$ & $-0.136$ & 608 \\
1.00 & $10^{-3}$ & $-0.452$ & 2027 \\
\bottomrule
\end{tabular}
\caption{Predicted $\delta K$ and Fe~K$\alpha$ centroid shift
for Cyg~X-1 as a function of Eddington ratio,
at $r=6M$, $\theta=\pi/2$, $a=0.98M$.
At $\theta=\pi/2$: $\Sigma=r^2=36M^2$ (since $\cos(\pi/2)=0$).
XRISM/Resolve energy resolution: 7\,eV.}
\label{tab:lum}
\end{table}

\subsection{Signature~III: Inclination dependence}

From equation~\eqref{eq:dust}, the ratio of $|\delta K|$ at
the pole to the equator is
\begin{equation}
  \frac{|\delta K|_{\theta=0}}{|\delta K|_{\theta=\pi/2}}
  = \frac{r^2+a^2}{r^2} = 1 + \left(\frac{a}{r}\right)^2.
\end{equation}
For Cyg~X-1 at $r=6M$: ratio $= 1+(0.98/6)^2\approx1.027$.

This 2.7\% inclination asymmetry depends only on $a/r$ and is
independent of $\rho_0$.
Comparing systems with similar $a$ but different viewing
inclinations can isolate this signature without knowing the
absolute accretion density.

%===================================================================
\section{Numerical predictions}\label{sec:pred}

Table~\ref{tab:targets} summarises $\delta K$ predictions for
four well-studied XRISM targets.

\begin{table}[h]
\centering
\renewcommand{\arraystretch}{1.3}
\begin{tabular}{@{}lcccccc@{}}
\toprule
System & $a/M$ & $\rho_0\,[M^{-2}]$ & $\delta K$ & $\delta E$ & SNR \\
\midrule
Cyg~X-1   & 0.98 & $10^{-3}$ & $-0.46$ & 2.1\,keV & $300\sigma$ \\
GX~339-4  & 0.90 & $5\times10^{-4}$ & $-0.23$ & 1.0\,keV & $148\sigma$ \\
M87*      & 0.90 & $10^{-4}$ & $-0.046$ & 207\,eV & $30\sigma$ \\
Sgr~A*    & 0.50 & $10^{-5}$ & $-0.005$ & 20\,eV & $3\sigma$ \\
\bottomrule
\end{tabular}
\caption{Predicted $K$-deviation and Fe~K$\alpha$ centroid shift
for four XRISM targets, evaluated at $r=6M$, $\theta=\pi/2$
($\Sigma=r^2$ since $\cos(\pi/2)=0$).
SNR is relative to XRISM/Resolve resolution of 7\,eV.
Accretion densities from GRMHD simulation estimates.}
\label{tab:targets}
\end{table}

Cyg~X-1 and GX~339-4 offer the highest predicted signal.
Both have extensive XRISM/Resolve observations in the public
archive (DARTS).

%===================================================================
\section{Discussion}\label{sec:disc}

\subsection{Separation from spin systematics}

Current Fe~K$\alpha$ spin measurements with {\tt RELXILL}
assume pure Kerr vacuum.
If $\delta K\neq0$, the best-fit spin parameter $a$ absorbs
part of the $\delta K$ signal, introducing a systematic error.
The direction of the bias is: matter redshifts the line,
which mimics a \emph{lower} spin (since lower spin also
redshifts the line at fixed inclination).

The luminosity correlation~\eqref{eq:correlation} breaks this
degeneracy: spin does not vary with luminosity, but $\delta K$
does.
A multi-epoch analysis of Cyg~X-1 at different luminosity
states, testing whether the iron line centroid correlates
linearly with $L_X$, provides a direct test of $\delta K$.

\subsection{Limitations}

The prediction~\eqref{eq:deltaK} is exact at the level of the
Raychaudhuri equation and the Einstein equations.
The numerical estimates in Tables~\ref{tab:lum}--\ref{tab:targets}
use GRMHD-motivated density values and assume
$u_\mu\ell^\mu\approx1$; detailed GRMHD modeling would
refine these estimates.
The iron line centroid formula~\eqref{eq:deltaE} is a
leading-order approximation; a full ray-tracing calculation
in the modified spacetime is deferred to future work.

\subsection{Relation to existing work}

The $\delta K$ formula~\eqref{eq:deltaK_full} is derived
from the self-consistency Proposition of~\cite{K1note};
it does not require any modification of GR.
The effect is distinct from non-Kerr deviations studied in
the bumpy black hole literature~\cite{Collins2004}:
here the spacetime satisfies the full Einstein equations,
and $\delta K$ measures the departure from \emph{vacuum}
rather than from Kerr geometry.

%===================================================================
\section{Conclusion}\label{sec:conc}

We have derived the exact formula
\begin{equation*}
  \delta K = -4\pi\Sigma(\rho+p)(u_\mu\ell^\mu)^2
\end{equation*}
for the departure of the $K$-functional from unity in an
accreting Kerr spacetime, and identified three observable
signatures in the Fe~K$\alpha$ line accessible to
XRISM/Resolve.
The luminosity correlation $\delta E(t)\propto L_X(t)$ is
the cleanest test: it is absent in standard Kerr reflection
and predicted to be ~2\,keV per Eddington luminosity for Cyg~X-1.

We invite observational collaborators with access to XRISM
data and {\tt RELXILL} expertise to test these predictions.
Confirmation would constitute the first observational evidence
for the $K=1$ self-consistency framework as a vacuum diagnostic
in astrophysical black holes.

\subsection*{Acknowledgements}
No funding was received for this work.

\subsection*{Data availability}
No new data were generated. XRISM data are publicly available
at \url{https://darts.isas.jaxa.jp/astro/xrism/}.

\begin{thebibliography}{99}

\bibitem{K1note}
Y.\,Y.\,N.~Li,
Self-consistency of the Kerr principal null congruence:
an exact identity and its equivalence with the vacuum condition.
Submitted to Class.\ Quantum Grav.\ (2026).

\bibitem{Fabian2000}
A.\,C.~Fabian et al.,
Broad iron lines in active galactic nuclei.
Publ.\ Astron.\ Soc.\ Pac.\ \textbf{112}, 1145 (2000).
\url{https://doi.org/10.1086/316610}

\bibitem{Reynolds2014}
C.\,S.~Reynolds,
Measuring black hole spin using X-ray reflection spectroscopy.
Space Sci.\ Rev.\ \textbf{183}, 277 (2014).
\url{https://doi.org/10.1007/s11214-013-0006-6}

\bibitem{Garcia2014}
J.~Garc\'ia et al.,
Improved reflection models of black-hole accretion disks.
Astrophys.\ J.\ \textbf{782}, 76 (2014).
\url{https://doi.org/10.1088/0004-637X/782/2/76}

\bibitem{Collins2004}
N.\,A.~Collins, S.\,A.~Hughes,
Towards a formalism for mapping the spacetimes of massive
compact objects.
Phys.\ Rev.\ D \textbf{69}, 124022 (2004).
\url{https://doi.org/10.1103/PhysRevD.69.124022}

\end{thebibliography}

\end{document}
